% !TEX root = ../../main.tex

\section{可行性分析}

\subsection{时间可行性}

本课题的开发与实现主要安排在 2026 年 3 月至 2026 年 5 月期间，整体时间分配较为合理，能够覆盖需求分析、系统设计、功能实现、联调测试和论文撰写等主要阶段。
课题实施过程中，可按照“移动端训练功能实现—服务端接口与数据管理实现—手势识别服务联调—后台管理端完善—句子视频识别实验—系统测试与论文整理”的顺序逐步推进，使系统开发过程具有较清晰的阶段性目标。

从任务复杂度来看，虽然本系统涉及 HarmonyOS 移动端、Spring Boot 服务端、Python 手势识别服务以及 Web 管理端等多个部分，但各模块之间职责相对明确，可以采用分模块开发、逐步联调的方式降低开发风险。
其中，移动端主要负责课程学习、数字人展示、训练交互与识别结果展示；
服务端主要负责业务数据管理、文件资源管理、识别服务调用和语义修正结果组织；
识别服务主要负责关键点提取、特征处理、实时识别和句子视频识别；
管理端则负责手势资源、样本数据和模型版本维护。
这样的分工有利于在有限周期内逐步完成系统实现。

此外，本课题并未将目标设定为完整的开放域连续手语翻译系统，而是聚焦于“数字人动作示范 + 固定训练项手势识别 + 有限词表句子视频识别 + 后台资源与模型管理”的实现路径，避免了在毕业设计周期内引入过于复杂的大规模语料构建与端到端翻译模型训练任务。
结合目前已经完成的系统原型、功能模块和联调基础，可以认为本课题在时间安排上具有较好的可行性。

\subsection{技术可行性}

本系统采用的关键技术均具有较成熟的工程实践基础，且与本科阶段所学习的移动开发、Web 开发、数据库系统、软件工程和人工智能应用等知识联系紧密，具备较好的技术可行性。

在移动端方面，系统基于 HarmonyOS 平台，采用 ArkTS 和 ArkUI 实现课程浏览、训练交互、数字人资源展示、训练结果反馈和句子识别结果展示等功能。
HarmonyOS 提供了较完善的页面构建、状态管理和多媒体能力支持，能够满足本系统在移动端训练场景下的功能需求\cite{harmonyos-multi-device-deployment}。
在服务端方面，系统基于 Spring Boot 构建业务接口，并结合 MyBatis、MySQL、Redis 与 MinIO 实现数据存储、缓存管理和媒体资源管理，这些技术在 Web 应用开发中应用广泛，文档丰富、生态完善，能够较稳定地支撑系统业务实现。

在手势识别方面，系统采用 Python 作为识别服务实现语言，并结合 MediaPipe 与 TensorFlow 完成关键点提取、特征构造、模型训练和实时识别。
与直接在移动端部署复杂识别模型相比，本系统采用“移动端采集图像帧或上传视频—服务端转发或编排—Python 识别服务处理—结果返回客户端”的协同方案，更符合当前课题条件下的工程实现能力。
该方案既能够利用 MediaPipe 对手部和姿态关键点进行提取，又能通过独立识别服务降低移动端负担，提高模型维护与更新的灵活性\cite{lugaresi-mediapipe-2019,zhang-mediapipe-hands-2020}。

在句子视频识别方面，系统将任务限定在有限词表和受限句子场景下，通过动作窗口分析、词级候选序列生成和语义修正完成初步链路验证。
该方案不依赖大规模连续手语语料，也不直接承诺开放域翻译能力，而是将视觉识别结果和自然语言后处理进行组合，符合毕业设计阶段的实验边界。
与此同时，后台管理端基于 Vue 3 和 Element Plus 实现手势分类、手势资源、样本数据、训练管理和模型版本管理等功能，形成了较完整的系统支撑环境。
综上，本系统所采用的技术路线成熟、实现路径清晰，具有较强的技术可行性。

\subsection{操作可行性}

本系统面向的主要用户包括移动端训练用户以及后台管理端维护人员，不同角色的使用目标较为明确，系统交互流程也相对清晰，因此具有较好的操作可行性。

对于移动端用户而言，系统主要围绕课程学习、动作示范、训练识别和结果反馈展开，用户只需按照“浏览课程—进入训练项—观看数字人动作示范—开始练习—查看识别结果与反馈”的基本路径即可完成使用。
在句子视频识别场景中，用户可以上传一段手语句子视频，系统返回识别词序列、中文展示和语义修正后的句子结果。
该流程符合一般移动应用的使用习惯，页面功能入口明确，操作步骤较少，便于用户理解和上手。

对于后台管理端用户而言，系统采用可视化界面组织手势分类管理、手势资源管理、样本管理、训练管理和模型版本管理等功能，便于管理员对课程内容、训练资源和识别模型进行统一维护。
管理端用户无需直接接触底层代码或模型文件处理细节，只需通过页面化操作即可完成资源维护、样本查看、训练任务执行和模型版本发布等工作，这有助于降低系统维护门槛，提高后台操作效率。

此外，系统在整体设计中将普通用户训练流程与模型维护流程相对分离：移动端用户主要关注课程学习和训练体验，后台维护人员主要关注资源配置、样本管理与模型发布。
这样的角色分工有利于减少不必要的复杂操作，避免不同功能相互干扰，从而提高系统整体的可用性和可操作性。
